\documentclass[11pt,a4paper]{article}
\usepackage[utf8]{inputenc}
\usepackage{amsmath,amssymb}
\usepackage[margin=2.5cm]{geometry}
\usepackage{booktabs}

\newcommand{\eps}{\varepsilon}
\newcommand{\teps}{\underset{\sim}{\varepsilon}}
\newcommand{\rot}{\operatorname{rot}}
\newcommand{\Inc}{\operatorname{Inc}}

\title{From the Strain Equation to the Compatibility Equation}
\author{}
\date{}

\begin{document}
\maketitle

\section*{The idea: eliminate the displacement}

The strain equation gives $6$ components $\eps_{ij}$ from only $3$ unknowns $u_i$.
Differentiating enough times lets us combine the $\eps$'s so that every $u$ term cancels.
What remains is a relation among the $\eps$'s alone: the compatibility equation.
Throughout, a comma denotes partial differentiation, repeated indices are summed, and $\boldsymbol u$ is assumed of class $C^3$ so that derivatives commute (Schwarz's theorem).

\section{The 2D case}

In 2D, three strain components depend on two displacements $u$ and $v$:
\[
  \eps_{xx}=\frac{\partial u}{\partial x},\qquad
  \eps_{yy}=\frac{\partial v}{\partial y},\qquad
  2\eps_{xy}=\frac{\partial u}{\partial y}+\frac{\partial v}{\partial x}.
\]
Differentiate each equation so that every term contains the same third-order derivatives of $u$ and $v$:
\[
  \frac{\partial^2\eps_{xx}}{\partial y^2}=\frac{\partial^3 u}{\partial x\,\partial y^2},\qquad
  \frac{\partial^2\eps_{yy}}{\partial x^2}=\frac{\partial^3 v}{\partial y\,\partial x^2},\qquad
  2\frac{\partial^2\eps_{xy}}{\partial x\,\partial y}=\frac{\partial^3 u}{\partial x\,\partial y^2}+\frac{\partial^3 v}{\partial x^2\,\partial y}.
\]
The right-hand side of the third equation is the sum of the first two, hence
\[
  \boxed{\frac{\partial^2\eps_{xx}}{\partial y^2}+\frac{\partial^2\eps_{yy}}{\partial x^2}=2\frac{\partial^2\eps_{xy}}{\partial x\,\partial y}.}
\]

\section{The general 3D case}

\paragraph{Step 1 -- differentiate the strain twice.}
\[
  \eps_{ij}=\tfrac12\left(u_{i,j}+u_{j,i}\right)
  \quad\Longrightarrow\quad
  \eps_{ij,kl}=\tfrac12\left(u_{i,jkl}+u_{j,ikl}\right).
\]

\paragraph{Step 2 -- express the four terms of the compatibility equation.}
\begin{center}
\begin{tabular}{@{}ll@{}}
\toprule
Term & Value in terms of $\boldsymbol u$\\
\midrule
$\eps_{ij,kk}$ & $\tfrac12 u_{i,jkk}+\tfrac12 u_{j,ikk}$\\[2pt]
$\eps_{kk,ij}$ & $u_{k,kij}$\\[2pt]
$\eps_{jk,ik}$ & $\tfrac12 u_{j,kik}+\tfrac12 u_{k,jik}$\\[2pt]
$\eps_{ik,jk}$ & $\tfrac12 u_{i,kjk}+\tfrac12 u_{k,ijk}$\\
\bottomrule
\end{tabular}
\end{center}

\paragraph{Step 3 -- add the first two, subtract the last two.}
Since derivative order is irrelevant ($u_{j,kik}=u_{j,ikk}$, $u_{k,jik}=u_{k,kij}$, \dots), the terms cancel in pairs:
\begin{itemize}
  \item $\tfrac12 u_{i,jkk}$ cancels with $\tfrac12 u_{i,kjk}$;
  \item $\tfrac12 u_{j,ikk}$ cancels with $\tfrac12 u_{j,kik}$;
  \item $u_{k,kij}$ cancels with $\tfrac12 u_{k,jik}+\tfrac12 u_{k,ijk}$.
\end{itemize}
Therefore
\[
  \boxed{\;\frac{\partial^2\eps_{ij}}{\partial x_k\partial x_k}+\frac{\partial^2\eps_{kk}}{\partial x_i\partial x_j}
  -\frac{\partial^2\eps_{jk}}{\partial x_i\partial x_k}-\frac{\partial^2\eps_{ik}}{\partial x_j\partial x_k}=0.\;}
\]

\section{A more systematic route: via the rotation}

This route explains \emph{where} the combination comes from, instead of checking it after the fact.
Split the displacement gradient into its symmetric and antisymmetric parts:
\[
  u_{i,j}=\eps_{ij}+\omega_{ij},\qquad \omega_{ij}=\tfrac12\left(u_{i,j}-u_{j,i}\right),
\]
where $\omega$ is the infinitesimal rotation. One differentiation gives its gradient from $\teps$ alone:
\[
  \omega_{ij,k}=\tfrac12\left(u_{i,jk}-u_{j,ik}\right)=\eps_{ik,j}-\eps_{jk,i},
\]
since $\eps_{ik,j}-\eps_{jk,i}=\tfrac12\left(u_{i,kj}+u_{k,ij}-u_{j,ki}-u_{k,ji}\right)=\tfrac12\left(u_{i,jk}-u_{j,ik}\right)$.

For $\omega$ to be a genuine function, its second derivatives must commute: $\omega_{ij,kl}=\omega_{ij,lk}$. Now
\[
  \omega_{ij,kl}=\eps_{ik,jl}-\eps_{jk,il},\qquad
  \omega_{ij,lk}=\eps_{il,jk}-\eps_{jl,ik},
\]
so the condition reads
\begin{equation}\label{eq:full}
  \eps_{ik,jl}+\eps_{jl,ik}-\eps_{il,jk}-\eps_{jk,il}=0,
\end{equation}
which is the full Saint-Venant system (81 scalar equations, 6 of them independent in 3D).

\paragraph{Recovering the contracted form.}
Simply setting $l=k$ in~\eqref{eq:full} gives the trivial identity $0=0$. Instead, take the component with indices $(i,k,j,k)$, that is $j\to k$, $k\to j$, $l\to k$, and sum over $k$:
\[
  \eps_{ij,kk}+\eps_{kk,ij}-\eps_{ik,jk}-\eps_{jk,ik}=0,
\]
which is the boxed equation of Section~2. In tensor notation the same condition reads
\[
  \Inc(\teps)=\rot\rot(\teps)=0.
\]

\paragraph{Why the condition is also sufficient.}
If~\eqref{eq:full} holds on a simply connected domain, the relation $\omega_{ij,k}=\eps_{ik,j}-\eps_{jk,i}$ can be integrated to obtain $\omega$; then $u_{i,j}=\eps_{ij}+\omega_{ij}$ can be integrated to obtain $\boldsymbol u$. The resulting displacement reproduces $\teps$, and is unique up to a rigid-body motion.

\end{document}
